\documentclass[aspectratio=169,10pt]{beamer}
\usepackage[spanish]{babel}
\usepackage{fontspec}
\setsansfont[Path=../assets/fonts/,UprightFont=InterVariable.ttf,BoldFont=InterVariable.ttf,ItalicFont=InterVariable.ttf,BoldItalicFont=InterVariable.ttf]{InterVariable.ttf}
\setmainfont[Path=../assets/fonts/]{InterVariable.ttf}
\usefonttheme{professionalfonts}
\usepackage{amsmath,amssymb,tikz,booktabs}
\usetikzlibrary{arrows.meta,positioning,calc}
\definecolor{USMBlue}{HTML}{003876}
\definecolor{USMDeep}{HTML}{002A5C}
\definecolor{USMCyan}{HTML}{00AEEF}
\definecolor{USMYellow}{HTML}{FDB913}
\definecolor{USMPale}{HTML}{EAF5FB}
\definecolor{Ink}{HTML}{12233B}
\definecolor{SoftGray}{HTML}{F2F5F8}
\definecolor{Warn}{HTML}{AA5B00}
\setbeamercolor{normal text}{fg=Ink,bg=white}
\setbeamercolor{frametitle}{fg=white,bg=USMBlue}
\setbeamercolor{structure}{fg=USMBlue}
\setbeamercolor{block title}{fg=white,bg=USMBlue}
\setbeamercolor{block body}{fg=Ink,bg=USMPale}
\setbeamercolor{alerted text}{fg=USMBlue}
\setbeamertemplate{navigation symbols}{}
\setbeamertemplate{itemize items}[circle]
\setbeamertemplate{footline}{
 \leavevmode\hbox{
 \begin{beamercolorbox}[wd=.88\paperwidth,ht=2.6ex,dp=.8ex,leftskip=1em]{footline}
 \color{USMBlue}\scriptsize UNIVERSIDAD TÉCNICA FEDERICO SANTA MARÍA
 \end{beamercolorbox}
 \begin{beamercolorbox}[wd=.12\paperwidth,ht=2.6ex,dp=.8ex,rightskip=1em plus1fil]{footline}
 \color{USMBlue}\scriptsize\hfill\insertframenumber
 \end{beamercolorbox}}}
\tikzset{
 stagebox/.style={draw=USMBlue,line width=.8pt,rounded corners=3pt,fill=USMPale,
  align=center,inner sep=5pt,minimum height=1.05cm,font=\small},
 active/.style={stagebox,fill=USMBlue,text=white},
 risk/.style={stagebox,fill=USMYellow!22,draw=Warn},
 arr/.style={-{Latex[length=2.5mm]},line width=1pt,draw=USMCyan},
 weak/.style={-{Latex[length=2.5mm]},line width=.8pt,dashed,draw=Warn}
}
\newcommand{\key}[1]{{\color{USMBlue}\bfseries #1}}
\newcommand{\captionline}[1]{\vspace{.18cm}{\color{USMBlue}\small #1}}
\begin{document}
{
\setbeamercolor{background canvas}{bg=USMBlue}
\begin{frame}[plain]
\vspace{.95cm}
{\color{USMCyan}\large CONTROL Y SEGURIDAD NUCLEAR}\\[.55cm]
{\color{white}\Huge\bfseries Control predictivo\\[.15cm] informado por riesgo}\\[.45cm]
{\color{white}\Large Contexto, resultados y vías de exploración}\\[.8cm]
{\color{USMYellow}\rule{9.2cm}{2pt}}\\[.25cm]
{\color{white}\normalsize DT--SMDFR \quad|\quad Universidad Técnica Federico Santa María}\\[.12cm]
{\color{white}\small Octubre de 2026}
\end{frame}
}

\begin{frame}{Pregunta central}
\vspace{.2cm}
\begin{block}{Decisión de control}
¿Qué acción reduce el riesgo térmico futuro, mantiene el seguimiento y puede calcularse dentro del tiempo disponible?
\end{block}
\vspace{.35cm}
\begin{center}
\begin{tikzpicture}[node distance=5mm]
\node[stagebox,minimum width=2.15cm] (o) {Estado\\observado};
\node[stagebox,minimum width=2.15cm,right=of o] (d) {Acción de\\control};
\node[stagebox,minimum width=2.15cm,right=of d] (f) {Futuros\\posibles};
\node[risk,minimum width=2.15cm,right=of f] (r) {Costo y\\riesgo};
\node[active,minimum width=2.15cm,right=of r] (a) {Mejor acción\\factible};
\draw[arr] (o)--(d);
\draw[arr] (d)--(f);
\draw[arr] (f)--(r);
\draw[arr] (r)--(a);
\end{tikzpicture}
\end{center}
\captionline{Meta: usar el riesgo calculado para comparar cada acción candidata.}
\end{frame}

\begin{frame}{Contexto actual: riesgo después de la decisión}
\vspace{.15cm}
\begin{center}
\begin{tikzpicture}[node distance=6mm]
\node[stagebox,minimum width=2.55cm] (t) {Gemelo digital\\estado y pronóstico};
\node[active,minimum width=2.65cm,right=of t] (c) {Control predictivo\\seguimiento + margen};
\node[stagebox,minimum width=2.5cm,right=of c] (u) {Acción del\\ACS};
\node[risk,minimum width=2.55cm,below=1.05cm of c] (p) {Árbol de eventos\\PSM a posteriori};
\draw[arr] (t)--(c);
\draw[arr] (c)--(u);
\draw[arr] (u.south) |- (p.east);
\draw[weak] (p.west) -- ++(-1.6,0) node[left,font=\small,align=right] {no retroalimenta\\esta búsqueda};
\end{tikzpicture}
\end{center}
\captionline{ACS: enfriamiento auxiliar. PSM: margen probabilístico de seguridad.}
\vspace{.05cm}
\begin{columns}[T]
\column{.48\textwidth}
\begin{block}{Lo que ya hace}
Predice temperatura y evalúa seguimiento, cambios de caudal y una penalización térmica determinista.
\end{block}
\column{.48\textwidth}
\begin{block}{Lo que aún falta}
Ponderar rutas por probabilidad y usar el riesgo probabilístico para elegir cada acción candidata.
\end{block}
\end{columns}
\end{frame}

\begin{frame}{Cuatro desafíos que definen el problema}
\begin{columns}[T]
\column{.49\textwidth}
\begin{block}{1. Objetivo comparable}
El costo actual concatena nodos de un árbol; las rutas no se ponderan y el número de muestras puede cambiar.
\end{block}
\vspace{.1cm}
\begin{block}{2. Riesgo dentro del control}
El margen determinista mide una trayectoria; falta probabilidad de falla y cola térmica.
\end{block}
\column{.49\textwidth}
\begin{block}{3. Calidad de la solución}
Una buena búsqueda no prueba optimalidad. Hace falta una cota inferior válida y un candidato factible.
\end{block}
\vspace{.1cm}
\begin{block}{4. Protecciones múltiples}
Con dos sistemas activos importan disponibilidad, orden de activación, retardos y respuesta pasiva.
\end{block}
\end{columns}
\end{frame}

\begin{frame}{Resultado 1: costo local y diferencia con la referencia}
\begin{columns}[T]
\column{.52\textwidth}
\begin{block}{Consulta fija, ejecución local}
\begin{tabular}{@{}lr@{}}
Seguimiento térmico & $0.348407$\\
Cambios de caudal & $0.013619$\\
Penalización de seguridad & $0$\\ \midrule
Costo total, candidato factible & $\mathbf{0.362026}$
\end{tabular}
\end{block}
\column{.44\textwidth}
\begin{block}{Lectura física}
Margen térmico mínimo pronosticado: \key{64.42 K} respecto de $1244$ K.
\end{block}
\begin{block}{Reproducibilidad}
El costo archivado es $0.326151$; la ejecución local difiere \key{+11.0\,\%}.
\end{block}
\end{columns}
\vspace{.15cm}
\captionline{La factibilidad se verificó en las muestras del piloto. El margen mostrado no es una garantía probabilística.}
\end{frame}

\begin{frame}{Resultado 2: efecto del margen determinista en los barridos}
\vspace{.2cm}
\begin{center}
\begin{tikzpicture}
\node[stagebox,minimum width=4.65cm,minimum height=2.15cm] (a) at (0,0)
 {\textbf{Convergencia}\\[.12cm]\Large 2 de 49\\[.08cm]\small configuraciones\\cambian la acción};
\node[stagebox,minimum width=4.65cm,minimum height=2.15cm] (b) at (5.35,0)
 {\textbf{Fase aguda}\\[.12cm]\Large 0 de 49\\[.08cm]\small configuraciones\\cambian la acción};

\end{tikzpicture}
\end{center}
\vspace{.12cm}
\begin{itemize}
\item En convergencia, los dos cambios requieren un umbral por debajo de la banda térmica alcanzable.
\item En fase aguda, el seguimiento ya lleva el enfriamiento al límite permitido.
\end{itemize}
\captionline{Resultado del barrido estudiado; no implica que una métrica probabilística tenga el mismo comportamiento.}
\end{frame}

\begin{frame}{Resultado 3: cotas de optimalidad en el ejemplo y el piloto}
\vspace{.15cm}
\begin{columns}[T]
\column{.49\textwidth}
\begin{block}{Ejemplo híbrido de dos modos}
\[
0.499024391\leq J^\star\leq 0.500000000
\]
Brecha $0.000975609$; óptimo analítico $0.5$.
\end{block}
\column{.49\textwidth}
\begin{block}{Piloto del reactor}
\[
0\leq J^\star\leq 0.362026
\]
La cota inferior disponible no mide cercanía al óptimo.
\end{block}
\end{columns}
\vspace{.3cm}
\begin{center}
\begin{tikzpicture}[x=1cm,y=1cm]
\node[stagebox,minimum width=2.4cm] (d) at (0,0) {Caja de\\controles};
\node[stagebox,minimum width=2.4cm] (e) at (3.1,0) {Todas las ramas\\y temperaturas};
\node[active,minimum width=2.4cm] (l) at (6.2,0) {Cota inferior\\de la caja};
\draw[arr] (d)--(e);
\draw[arr] (e)--(l);
\end{tikzpicture}
\end{center}
\captionline{La pieza pendiente es encerrar rigurosamente el pronóstico térmico para todos los controles y ramas.}
\end{frame}

\begin{frame}{Control informado por riesgo: evaluar cada candidato}
\vspace{.05cm}
\begin{center}
\resizebox{.94\linewidth}{!}{%
\begin{tikzpicture}[node distance=5mm]
\node[active,minimum width=2.0cm] (u) {Candidato\\de control};
\node[stagebox,minimum width=2.0cm,right=of u] (s) {Mismas\\realizaciones};
\node[stagebox,minimum width=2.1cm,right=of s] (r) {Rutas, temperaturas\\y pesos};
\node[risk,minimum width=2.0cm,right=of r] (m) {Costo esperado\\y riesgo};
\node[active,minimum width=1.9cm,right=of m] (o) {Comparar\\candidatos};
\draw[arr] (u)--(s);
\draw[arr] (s)--(r);
\draw[arr] (r)--(m);
\draw[arr] (m)--(o);
\end{tikzpicture}%
}
\end{center}
\vspace{.25cm}
\[
\min_{\pi}\ \mathbb E[J(\pi,\omega)]+\lambda_R R(\pi)
\qquad
\text{sujeto a}\quad
\mathbb P\!\left(\max_k T_k>1244\ {\rm K}\right)\leq\varepsilon .
\]
\vspace{.05cm}
\begin{columns}[T]
\column{.49\textwidth}
\begin{block}{Métrica principal a probar}
Probabilidad de exceder el límite durante un horizonte común.
\end{block}
\column{.49\textwidth}
\begin{block}{Métricas complementarias}
Margen por cuantil y costo medio de los peores casos (CVaR).
\end{block}
\end{columns}
\end{frame}

\begin{frame}{Vías de exploración más prometedoras}
\begin{columns}[T]
\column{.32\textwidth}
\begin{block}{1. Escenarios comunes}
Evaluar cada control sobre las mismas realizaciones. Comparar costo y riesgo sin que cambie el ruido de muestreo.
\end{block}
\column{.32\textwidth}
\begin{block}{2. Fallas raras}
Usar muestreo por importancia para la cola; registrar pesos, tamaño efectivo e incertidumbre.
\end{block}
\column{.32\textwidth}
\begin{block}{3. Cotas globales}
Construir encierros térmicos por regiones de control y ramas; medir una brecha de optimalidad válida.
\end{block}
\end{columns}
\vspace{.3cm}
\begin{center}
\begin{tikzpicture}[node distance=5mm]
\node[stagebox,minimum width=3.05cm] (a) {Primero: decisión\\sensible al riesgo};
\node[stagebox,minimum width=3.05cm,right=of a] (b) {Después: cola\\estadísticamente estable};
\node[stagebox,minimum width=3.05cm,right=of b] (c) {En paralelo: calidad\\de la optimización};
\draw[arr] (a)--(b);
\draw[arr] (b)--(c);
\end{tikzpicture}
\end{center}
\end{frame}

\begin{frame}{Con varios sistemas, la estrategia es una secuencia causal}
\vspace{.05cm}
\begin{center}
\begin{tikzpicture}[node distance=8mm and 9mm]
\node[stagebox,minimum width=2.4cm] (obs) {Estado y demanda\\observados};
\node[active,minimum width=2.4cm,right=of obs] (acs) {Activar ACS\\si corresponde};
\node[stagebox,minimum width=2.35cm,right=of acs] (out) {Observar\\éxito o falla};
\node[active,minimum width=2.45cm,right=of out] (b) {Decidir sobre\\sistema B};
\node[risk,minimum width=2.65cm,below=1.15cm of out] (pas) {Protección pasiva\\según demanda y estado};
\draw[arr] (obs)--(acs);
\draw[arr] (acs)--(out);
\draw[arr] (out)--(b);
\draw[weak] (obs.south) |- (pas.west);
\draw[weak] (out.south)--(pas.north);
\end{tikzpicture}
\end{center}
\vspace{.2cm}
\begin{itemize}
\item ACS$\rightarrow$B y B$\rightarrow$ACS pueden producir riesgos y respuestas térmicas distintos.
\item La acción en cada instante usa sólo fallas y demandas ya observadas.
\item La protección pasiva responde a demanda y disponibilidad; el controlador decide sobre los sistemas activos.
\end{itemize}
\end{frame}

\begin{frame}{Secuencia de trabajo y criterio de éxito}
\begin{center}
\begin{tikzpicture}[node distance=7mm]
\node[stagebox,minimum width=2.5cm] (a) {1. Conciliar\\línea base};
\node[stagebox,minimum width=2.5cm,right=of a] (b) {2. Riesgo en\\cada candidato};
\node[stagebox,minimum width=2.5cm,right=of b] (c) {3. Validar fuera\\de muestra};
\node[stagebox,minimum width=2.5cm,right=of c] (d) {4. Dos sistemas\\y pasivo};
\draw[arr] (a)--(b);
\draw[arr] (b)--(c);
\draw[arr] (c)--(d);
\end{tikzpicture}
\end{center}
\vspace{.35cm}
\begin{block}{Una mejora será convincente cuando}
La nueva métrica cambie alguna decisión, reduzca el riesgo fuera de muestra a un costo térmico aceptable y cumpla el tiempo disponible por consulta.
\end{block}
\vspace{.12cm}
\captionline{En paralelo: obtener cotas para el modelo real y cuantificar la brecha de optimalidad.}
\end{frame}
\end{document}
